\chapter{Arquitectura híbrida residual LSTM-XGBoost para demanda diaria agregada}\label{cap:metodologia}

Este capítulo describe la arquitectura de modelado propuesta y el protocolo experimental
que la sustenta, sobre la matriz de 161 variables construida en el capítulo
\ref{cap:datos_features}. La propuesta central es un híbrido residual secuencial de dos
etapas: una red LSTM univariante (etapa 1) que modela la dinámica temporal de la propia
demanda, y un modelo XGBoost (etapa 2) que corrige, a partir de la matriz completa de
variables predictoras (calendario y meteorología retardada, pero también retardos y
estadísticos móviles de la propia demanda), el residuo que la etapa 1 deja sin explicar.
Formalmente, si
$\hat y^{\text{LSTM}}_t$ denota la predicción de la etapa 1 en el instante $t$, el residuo
que alimenta la etapa 2 se define como
\[
  e_t = y_t - \hat y^{\text{LSTM}}_t,
\]
y la predicción final del sistema combina ambas etapas por suma directa,
\[
  \hat y^{\text{final}}_t = \hat y^{\text{LSTM}}_t + \hat e^{\text{XGBoost}}_t .
\]

Conviene precisar en qué difiere este diseño del esquema híbrido ARIMA-red neuronal de
\textcite{zhang2003hybrid}, que es su referente más directo en la literatura y al que
se volverá al interpretar el resultado central del capítulo \ref{cap:resultados}
(\ref{sec:resultado_negativo_teoria}). En primer lugar, el orden de las etapas está
invertido: aquí es la red neuronal la que modela primero la componente no lineal de la
serie y un modelo basado en árboles el que corrige después, no al revés. En segundo
lugar, la etapa 2 de este TFM no modela el residuo como una serie temporal autorregresiva
sobre sus propios retardos, sino como una función de 161 variables predictoras (retardos
de la demanda, calendario, meteorología retardada, términos de Fourier) construidas de
forma independiente de la
dinámica del residuo mismo. Y en tercer lugar, y de manera no negociable para la validez
del diseño, las predicciones de la etapa 1 que se emplean para calcular ese residuo de
entrenamiento se generan siempre fuera de muestra (\emph{out-of-fold}, OOF), mediante un
esquema de validación cruzada expansivo hacia adelante (\ref{sec:esquema_oof}) y nunca a
partir de un único ajuste \emph{in-sample}: un residuo calculado sobre predicciones que el
propio modelo ya ha visto en entrenamiento está optimistamente sesgado hacia cero y es
estructuralmente distinto del residuo que la etapa 2 enfrentará realmente en inferencia,
de modo que aprender a corregirlo no transferiría a producción.

\section{Diagnóstico previo: ¿tiene el residuo estructura de calendario?}\label{sec:diagnostico_previo}

Antes de comprometerse con la arquitectura de dos etapas descrita más arriba conviene
comprobar, con evidencia y no por intuición, que existe una estructura de error
sistemática que una segunda etapa no lineal pueda efectivamente corregir: si un modelo
lineal ya captura todo el error explicable, ni el híbrido residual ni el repesado por
calendario de la sección \ref{sec:experimento_a} tendrían ningún margen de mejora
posible. Este diagnóstico exploratorio se calcula sobre el error absoluto de SARIMAX
(el baseline estadístico de la sección \ref{sec:baselines}, disponible ya en esta fase
del proyecto) agregado por régimen de calendario sobre las 1.282 observaciones
posteriores al recorte de calentamiento de 28 días (sección \ref{sec:particiones_fuga}),
y es deliberadamente exploratorio: se calcula sobre el periodo completo train+val+test
y no sobre una partición fuera de muestra, porque los días festivos son demasiado raros
(48 en total) para dividirlos de forma fiable en tres particiones y seguir teniendo un
tamaño de grupo interpretable.

El resultado confirma la premisa de la arquitectura híbrida: el error absoluto medio de
SARIMAX es 492.092 en festivos (20,50\,\% de error porcentual medio) y 262.194 en
sábado, frente a 212.383 en laborables en general, y el sesgo con signo es
sistemáticamente negativo en festivos ($-145.250$), es decir, SARIMAX
\textbf{sobreestima} la demanda festiva de forma consistente sin saber cuantificar
correctamente su magnitud.

El contraste más nítido, sin embargo, aísla el efecto de los
días puente del efecto de fin de semana comparando únicamente dentro del régimen
laborable: los 54 días laborables de tipo puente presentan un error absoluto medio de
518.634 frente a 192.191 en los 819 laborables ordinarios, es decir,
\textbf{los días puente concentran 2,70 veces el error absoluto medio de un laborable
corriente}. Esta es exactamente la estructura no lineal y condicionada al calendario
(grande en magnitud, concentrada en un subconjunto pequeño y bien identificable de
días) que la etapa 2 de la arquitectura híbrida (sección \ref{sec:etapa2_xgboost}) está
diseñada para aprender a partir de variables exógenas, y la que motiva directamente,
más adelante, tanto el repesado de días irregulares del experimento A (sección
\ref{sec:experimento_a}) como la lectura del desglose de error por tipo de día del
capítulo \ref{cap:resultados} (sección \ref{sec:error_dia}). Que el mecanismo exista de
forma medible en este diagnóstico exploratorio, y que aun así el híbrido resultante no
supere a sus componentes de una sola etapa, es precisamente el contraste que sostiene
el hallazgo central de este TFM: el problema no es que no haya señal de calendario que
corregir, sino que la etapa 1 de la que parte esa corrección es demasiado débil
(sección \ref{sec:resultado_negativo_teoria}).

\section{Baselines}\label{sec:baselines}

Antes de evaluar cualquier arquitectura de aprendizaje automático es necesario fijar la
referencia frente a la que se mide: un conjunto de modelos de referencia sencillos que
establezcan cuánta señal puede explicarse sin recurrir a redes neuronales ni a variables
exógenas. El primero, la persistencia $\hat y_t = y_{t-1}$, resulta ser un modelo de referencia
particularmente débil sobre esta serie: como se documentó en la sección \ref{sec:eda}, la
demanda salta en torno a 2,5 millones de viajes diarios en cada transición entre régimen
laborable y régimen de fin de semana, una transición que ocurre dos veces por semana, y la
persistencia ignora por construcción cualquier estructura de calendario. El resultado es
un coeficiente de determinación negativo tanto en entrenamiento como en test, lo que
indica que ni siquiera predecir la media de la serie sería peor. Por ello, el baseline
ingenuo relevante en este trabajo es la persistencia estacional $\hat y_t = y_{t-7}$, que
sí incorpora el ciclo semanal al comparar cada día con el mismo día de la semana anterior;
superar a la persistencia estacional, y no a la persistencia simple, es el umbral
que cualquier modelo posterior debe rebasar para justificar su complejidad adicional.

El baseline estadístico fuerte de esta familia es un modelo SARIMAX, seleccionado
mediante una búsqueda en rejilla de 324 combinaciones de órdenes
$(p,d,q)\times(P,D,Q,s)$, con $p,q,P,Q \in \{0,1,2\}$, $d,D \in \{0,1\}$ y periodo
estacional fijo $s=7$, el ciclo semanal identificado en el análisis exploratorio. La
selección se realiza exclusivamente por el criterio de información de Akaike (AIC)
calculado sobre el conjunto de entrenamiento, deliberadamente sin consultar en ningún
momento el error de validación: hacerlo convertiría la validación en un conjunto de
ajuste adicional y contaminaría la comparación que se plantea en el capítulo
\ref{cap:resultados}. El orden seleccionado, resumido en la tabla \ref{tab:sarimax_orden},
es $\text{SARIMAX}(2,1,2)\times(0,1,2,7)$, con un AIC de entrenamiento de 24.616,89; las
cinco combinaciones mejor puntuadas se sitúan dentro de 0,8 puntos de AIC entre sí, de
modo que el orden concreto elegido debe leerse como \emph{una} especificación razonable de
esta complejidad, no como una verdad estructural descubierta de forma inequívoca.

\begin{table}[htbp]
  \centering
  \small
  \begin{tabular}{llrr}
\toprule
\rowcolor{naranja}
\textbf{Orden (p,d,q)} & \textbf{Orden estacional (P,D,Q,s)} & \textbf{AIC} & \textbf{N.º variables exógenas} \\
\midrule
\rowcolor{filaclara} (2,1,2) & (0,1,2,7) & 24.616,89 & 5 \\
\bottomrule
\end{tabular}

  \caption{\headlinecolor{Orden SARIMAX seleccionado por AIC sobre 324 combinaciones.}}
  \label{tab:sarimax_orden}
\end{table}

El vector de variables exógenas se limita deliberadamente a cinco columnas de calendario
(los indicadores de fin de semana y de día puente, y las tres categorías reales de tipo
de festividad, excluyendo \texttt{feat\_holiday\_type\_none} como nivel de referencia para
evitar la trampa de la variable ficticia frente al intercepto de la regresión), sin
incorporar ningún retardo ni estadístico móvil de la propia demanda: el componente
$(p,d,q)\times(P,D,Q,s)$ del modelo ya captura la autorregresión y la diferenciación
estacional de la serie \parencite{box1970time}, y añadir retardos hechos a mano expresaría la
misma dependencia dos veces a través de dos mecanismos en competencia.

Los parámetros del modelo se estiman una única vez sobre el conjunto de entrenamiento; a partir de ahí, la
predicción sobre validación y test es un pronóstico \emph{walk-forward} de un paso
mediante el filtro de Kalman (\texttt{.filter()}, sin reestimación de parámetros) con
\texttt{get\_prediction(dynamic=False)}, de modo que cada predicción en el instante $t$
condiciona exclusivamente sobre las observaciones reales hasta $t-1$. La evaluación de
SARIMAX corresponde, por tanto, a un escenario de predicción \textbf{de un día hacia
adelante}, actualizado diariamente una vez conocida la demanda observada del día anterior;
no representa un pronóstico estático de todo el periodo de evaluación realizado desde su
fecha inicial. Se prefiere este
esquema a un único pronóstico estático a 390 días porque este último decaería hacia la
media de la serie y mediría, en la práctica, una distancia de extrapolación en lugar de la
calidad real del modelo; se verificó además que los residuos no presentan autocorrelación
de primer orden relevante (aproximadamente $-0{,}03$) y que la predicción sigue de cerca
el valor real del propio día ($\text{MAE}=163.627$) y no el del día anterior
($\text{MAE}=909.082$), descartando así un desfase temporal accidental en la
implementación. Con este protocolo, SARIMAX alcanza un MAPE de test del 3,85\,\% y un
$R^2$ de 0,9672 (tabla \ref{tab:comparacion_test}, sección \ref{sec:tabla_maestra}),
constituyéndose en un baseline exigente que cualquier arquitectura de aprendizaje
automático de este TFM debe superar para justificar su mayor coste de desarrollo. Es el
mismo horizonte de un día al que se evalúan todos los demás modelos del estudio: la matriz
de variables solo contiene retardos $\geq 1$ y la ventana de la etapa 1 abarca $[t-W, t)$,
de modo que la comparación de la tabla maestra es entre pronósticos a un día vista en todos
los casos.

\section{Etapa 1: LSTM}\label{sec:etapa1_lstm}

La primera etapa del híbrido es una red LSTM \parencite{hochreiter1997lstm}
\textbf{univariante por diseño}: recibe únicamente la ventana de las $W$ observaciones
pasadas de \texttt{total} y debe inferir a partir de esa dinámica (sin ningún indicador de
día de la semana, festividad o puente) toda la estructura de calendario que la sección
\ref{sec:eda} mostró dominante. Esta restricción no es una limitación accidental, sino la
pregunta que la etapa 1 está diseñada para responder: cuánta de esa estructura puede
recuperarse de la memoria temporal de la serie por sí sola, dejando a la etapa 2 la tarea
explícita de aportar la información de calendario y meteorología que la etapa 1
estructuralmente no puede ver.

La longitud de la ventana $W$ se fijó mediante una comparación empírica, no una búsqueda
exhaustiva, entre tres valores candidatos, $W \in \{14, 28, 60\}$, entrenando en cada caso
un modelo final completo y evaluándolo únicamente sobre la partición de validación; el
conjunto de test se mantuvo deliberadamente fuera de esta decisión, de modo que la
elección de $W$ no pudiera beneficiarse, ni siquiera indirectamente, de información sobre
la partición de evaluación final. Los resultados se recogen en la tabla
\ref{tab:ventana_lstm}: $W=28$ resulta ganador en las cuatro métricas de validación
simultáneamente, con una brecha de MAE de aproximadamente el 11\,\% frente a sus dos
vecinos, una diferencia lo bastante amplia como para no atribuirse al azar de una única
ejecución, aunque el propio diseño del experimento, una ejecución por valor de $W$, obliga
a presentarla como evidencia documentada y no como el resultado de una búsqueda
sistemática de hiperparámetros.

\begin{table}[htbp]
  \centering
  \small
  \setlength{\tabcolsep}{4pt}
  \resizebox{\textwidth}{!}{%
    \begin{tabular}{>{\centering\arraybackslash}p{0.8cm}>{\centering\arraybackslash}p{2.1cm}>{\centering\arraybackslash}p{2.0cm}>{\centering\arraybackslash}p{1.7cm}>{\centering\arraybackslash}p{1.9cm}>{\centering\arraybackslash}p{1.9cm}>{\centering\arraybackslash}p{2.0cm}>{\centering\arraybackslash}p{1.5cm}}
\toprule
\rowcolor{naranja}
\textbf{W} & \textbf{Secuencias train} & \textbf{Épocas ejecutadas} & \textbf{Mejor época} & \textbf{MAE (val)} & \textbf{RMSE (val)} & \textbf{MAPE (val, \%)} & \textbf{R² (val)} \\
\midrule
\rowcolor{filaclara} 14 & 875 & 154 & 144 & 459.813 & 721.182 & 13,28 & 0,7607 \\
28 & 861 & 164 & 154 & 411.142 & 655.479 & 11,81 & 0,8023 \\
\rowcolor{filaclara} 60 & 829 & 87 & 77 & 456.063 & 692.825 & 12,89 & 0,7792 \\
\bottomrule
\end{tabular}
%
  }
  \caption{\headlinecolor{Comparación de longitud de ventana W (evidencia, no búsqueda exhaustiva: una ejecución final por W, test nunca consultado en la elección).}}
  \label{tab:ventana_lstm}
\end{table}

Conviene señalar que $W$ se eligió observando precisamente la
partición de validación: por tanto, las métricas de validación de toda la familia
LSTM (\texttt{lstm\_alone}, \texttt{hybrid}, \texttt{hybrid\_weighted}) son levemente
optimistas respecto de lo que cabría esperar de una evaluación totalmente independiente,
mientras que sus métricas de test, nunca consultadas durante esta elección, constituyen
la evaluación estrictamente limpia y son las que deben citarse como comparación
principal frente a SARIMAX y XGBoost.

La arquitectura final, entrenada sobre $W=28$, es deliberadamente compacta:
una capa LSTM de 32 unidades con desactivación aleatoria (\emph{dropout}) de 0,2, seguida de una capa densa de
salida lineal de una única neurona. Se optimiza con Adam ($\text{lr}=10^{-3}$) bajo error
cuadrático medio, en lotes de 32 secuencias, con parada anticipada de paciencia 10 y
restauración de los mejores pesos observados (\texttt{restore\_best\_weights}). El corte
de la parada anticipada se decide sobre una cola interna de las propias secuencias de
\emph{entrenamiento} (las últimas 129, correspondientes al periodo 2025-02-27 a
2025-07-05) y nunca sobre la partición de validación oficial: emplear la validación
oficial en ese punto convertiría la parada anticipada en un mecanismo de ajuste sobre
validación, sesgando de forma optimista y silenciosa cada cifra de validación que se
reporte más adelante en la tabla maestra de comparación (\ref{sec:tabla_maestra}).

Entrenada de este modo sobre las 889 filas de entrenamiento posteriores al calentamiento,
la etapa 1 en solitario queda claramente por debajo de SARIMAX en ambas particiones de
evaluación. Y, más revelador todavía, no llega a superar ni siquiera a la persistencia
estacional cuando ambas se comparan sobre las mismas filas fuera de muestra generadas por
el esquema de la sección \ref{sec:esquema_oof} (MAE OOF de 690.460 frente a 461.087).
Esto no constituye un defecto de la implementación: los resultados obtenidos indican
que la dinámica temporal univariante, por sí sola, no basta
para igualar a un modelo que sí observa el calendario de forma explícita. El detalle
numérico completo de esta comparación, junto con el resto de modelos de una sola etapa,
se presenta en la tabla maestra del capítulo \ref{cap:resultados}
(\ref{sec:tabla_maestra}).

Dos decisiones de ajuste, verificadas empíricamente y no asumidas de antemano, merecen
constar. Primero, el número máximo de épocas se elevó de 200 a 500 tras observar que el
fold 3 del esquema OOF alcanzaba las 200 épocas sin que la parada anticipada hubiera
actuado todavía —es decir, el límite artificial de épocas, no la convergencia del modelo,
era la restricción activa—, lo que habría entrenado folds de tamaño distinto hasta
presupuestos de cómputo efectivamente distintos y los habría hecho incomparables entre sí;
al elevar el límite, el fold 3 llegó a 228 épocas y su MAE mejoró de 288.194 a 282.122.
Segundo, se probó y se descartó el barajado de los lotes de entrenamiento: aunque barajar
secuencias ya construidas no constituiría en sí misma una fuga de información —cada
secuencia es autocontenida y la cola de parada anticipada se separa cronológicamente antes
de cualquier barajado—, medido sobre el fold 3 el barajado empeoró el resultado (MAE OOF de
288.000 a 309.000 aproximadamente), por lo que se mantiene \texttt{shuffle=False} por
evidencia, no por precaución teórica.

\section{Esquema de validación fuera de muestra (OOF)}\label{sec:esquema_oof}

El esquema descrito en esta sección es, con diferencia, el componente metodológico más
crítico del proyecto: la decisión de diseño de la que depende la validez de toda la etapa
2. Como se anticipó en la introducción del capítulo, las predicciones de la etapa 1 que
alimentan el cálculo del residuo de entrenamiento de la etapa 2 deben generarse siempre
fuera de muestra. La alternativa sería ajustar un único modelo LSTM sobre todo el
entrenamiento y calcular el residuo directamente sobre esas mismas predicciones
\emph{in-sample}, y produciría un residuo artificialmente pequeño y con una distribución
distinta de la que la etapa 2 encontrará en producción, precisamente porque el modelo ya
ha memorizado, en algún grado, las observaciones sobre las que se evalúa; la etapa 2
aprendería entonces a corregir una distribución de error que no existe fuera de muestra,
lo que se traduciría en un desempeño peor del esperado en inferencia real. Este esquema de
validación fuera de muestra para series temporales sigue la práctica establecida en la
literatura de evaluación de modelos de pronóstico
\parencite{bergmeir2012cv,cerqueira2020evaluating}, adaptada aquí a la generación de
predicciones de \emph{stacking} \parencite{wolpert1992stacking} en dos etapas.

Concretamente, se emplea una validación cruzada expansiva hacia adelante
(\emph{walk-forward}) de cinco folds sobre el periodo de entrenamiento posterior al
calentamiento, 2023-01-29 a 2025-07-05 (889 filas). El fold $k$ entrena un modelo LSTM
completo sobre el bloque $[0, s_k)$ y genera predicciones únicamente sobre el bloque
inmediatamente siguiente, $[s_k, s_{k+1})$, que nunca ha visto durante su entrenamiento;
cada fold reajusta además su propio escalador de normalización exclusivamente sobre su
propia porción de entrenamiento, sin reutilizar en ningún caso un escalador ajustado con
información de un bloque posterior. La tabla \ref{tab:folds_oof} recoge las métricas de
cada uno de los cinco folds, y la figura \ref{fig:oof_folds} muestra sus predicciones
fuera de muestra frente al valor real.

\begin{table}[htbp]
  \centering
  \small
  \setlength{\tabcolsep}{4pt}
  \resizebox{\textwidth}{!}{%
    \begin{tabular}{>{\centering\arraybackslash}p{1.0cm}>{\centering\arraybackslash}p{2.1cm}>{\centering\arraybackslash}p{1.9cm}>{\centering\arraybackslash}p{2.2cm}>{\centering\arraybackslash}p{2.1cm}>{\centering\arraybackslash}p{1.9cm}>{\centering\arraybackslash}p{1.9cm}}
\toprule
\rowcolor{naranja}
\textbf{Fold} & \textbf{Épocas ejecutadas} & \textbf{Mejor época} & \textbf{Loss train (final)} & \textbf{Loss val (final)} & \textbf{MAE} & \textbf{RMSE} \\
\midrule
\rowcolor{filaclara} 1 & 15 & 5 & 0,1195 & 0,0788 & 1.461.905 & 1.682.893 \\
2 & 153 & 143 & 0,0271 & 0,0509 & 659.425 & 882.524 \\
\rowcolor{filaclara} 3 & 228 & 218 & 0,0207 & 0,0281 & 282.122 & 397.512 \\
4 & 131 & 121 & 0,0208 & 0,0098 & 548.081 & 839.318 \\
\rowcolor{filaclara} 5 & 81 & 71 & 0,0251 & 0,0331 & 506.148 & 718.174 \\
\bottomrule
\end{tabular}
%
  }
  \caption{\headlinecolor{Métricas por fold del esquema OOF \emph{walk-forward} expansivo (5 folds).}}
  \label{tab:folds_oof}
\end{table}

\begin{figure}[htbp]
  \centering
  \includegraphics[width=0.78\textwidth]{fig_oof_folds.png}
  \caption{\headlinecolor{Predicción LSTM fuera de muestra (OOF) por fold del \emph{walk-forward} expansivo.}}
  \label{fig:oof_folds}
\end{figure}

Las 200 primeras filas de entrenamiento (2023-01-29 a 2023-08-16) no reciben ninguna
predicción OOF, puesto que evaluarlas exigiría un modelo entrenado con menos observaciones
que la ventana mínima inicial que el esquema considera aceptable; estas filas se
conservan en el fichero de salida marcadas explícitamente como \texttt{has\_oof=False},
con su motivo de exclusión indicado, en lugar de eliminarse sin más. La cobertura efectiva
resultante es de 689 de las 889 filas de entrenamiento (77,5\,\%), con un MAE OOF agregado
de 690.460 y un RMSE de 997.520.

El fold 1 merece un análisis detallado porque su comportamiento es \textbf{degenerado} y
la decisión que se toma sobre él condiciona directamente el conjunto de entrenamiento de
la etapa 2. Entrenado sobre apenas 200 observaciones (172 secuencias, reducidas a 147 tras
separar la cola de parada anticipada), sus predicciones fuera de muestra presentan una
desviación típica de solo el 0,199 de la de los valores reales y una correlación de
apenas 0,300 con ellos: el modelo, con tan poco historial disponible, nunca llega a
escapar de predecir un valor próximo a una constante, y la parada anticipada se activa ya
en la época 15 sobre una señal de validación interna construida con solo 25 secuencias,
demasiado ruidosa para guiar el entrenamiento de forma fiable. La figura
\ref{fig:oof_fold1_degenerado} contrasta visualmente esta degeneración con el fold 3, sano
por comparación (correlaciones entre 0,79 y 0,95 en los folds 2 a 5, frente al 0,300 del
fold 1).

\begin{figure}[htbp]
  \centering
  \includegraphics[width=0.78\textwidth,keepaspectratio]{fig_oof_fold1_degenerado.png}
  \caption{\headlinecolor{Fold 1 del esquema OOF es degenerado (predicción casi constante) frente al fold 3 (sano, referencia).}}
  \label{fig:oof_fold1_degenerado}
\end{figure}

La consecuencia práctica es que las 137 filas del fold 1 se \textbf{excluyen} del
conjunto de entrenamiento de la etapa 2 (\ref{sec:etapa2_xgboost}): sus residuos están
dominados por el fallo estructural de un modelo con historial insuficiente, no por la
estructura de calendario que se pretende que XGBoost aprenda, e incorporarlos
introduciría ruido de una naturaleza distinta a la del resto del conjunto residual. Esta
exclusión no se decide por intuición, sino que se somete a una comprobación explícita de
cobertura de categorías: para cada nivel categórico presente en el fold 1 (tipo de día,
tipo de festividad, indicador de día puente) se verifica que ese mismo nivel reaparece con
observaciones suficientes en los folds 2 a 5 restantes. Ninguna categoría queda huérfana
tras la exclusión, de modo que eliminar el fold 1 retira redundancia de cobertura y no
información única e insustituible; de haberse detectado una categoría con cobertura
exclusiva en el fold 1, el procedimiento habría detenido la exclusión en lugar de
proceder de todos modos.

\section{Etapa 2: XGBoost sobre el residuo}\label{sec:etapa2_xgboost}

Tras retirar las 200 filas sin predicción OOF y las 137 filas del fold 1 degenerado, el
conjunto de entrenamiento de la etapa 2 queda formado por 552 filas, desde el
2024-01-01 hasta el 2025-07-05, con las 161 variables predictoras de la sección
\ref{sec:ingenieria_features} y sin ningún valor nulo. La variable objetivo de esta etapa
no es la demanda sino el residuo definido en la introducción del capítulo,
$e_t = y_t - \hat y^{\text{LSTM}}_t$, calculado exclusivamente a partir de las
predicciones fuera de muestra de la sección anterior; su signo se fija de modo que la
corrección de la etapa 2 se \emph{suma} a la predicción de la etapa 1. El residuo
presenta una media de 109.601 y una desviación típica de 726.669 sobre un rango de
$[-3.128.729,\ +1.899.895]$: una dispersión del mismo orden de magnitud que la propia
señal de demanda, lo que anticipa que corregirlo es, en sí mismo, un problema de
predicción exigente.

XGBoost \parencite{chen2016xgboost} se entrena sobre este conjunto residual con los
hiperparámetros de la tabla \ref{tab:hiperparametros}, seleccionados mediante una rejilla
de 54 combinaciones evaluada sobre un \emph{holdout} interno cronológico (el último
15\,\% de las fechas del propio conjunto de entrenamiento del modelo, nunca la partición
de validación oficial) y reajustados después sobre la totalidad de los datos disponibles
para ese modelo. La configuración ganadora combina árboles relativamente profundos
($\text{max\_depth}=5$) con una tasa de aprendizaje baja (0,01) y un número elevado de
estimadores (300), junto con una regularización moderada por tamaño mínimo de nodo
($\text{min\_child\_weight}=3$); esta combinación resulta coherente con un conjunto de
entrenamiento reducido (552 filas) y una señal de calendario dispersa y de bajo peso
relativo frente al ruido introducido por la propia etapa 1.

\begin{table}[htbp]
  \centering
  \small
  \setlength{\tabcolsep}{4pt}
  \resizebox{\textwidth}{!}{%
    \begin{tabular}{>{\raggedright\arraybackslash}p{3.0cm}>{\centering\arraybackslash}p{2.3cm}>{\centering\arraybackslash}p{1.8cm}>{\centering\arraybackslash}p{2.1cm}>{\centering\arraybackslash}p{2.0cm}>{\centering\arraybackslash}p{2.3cm}}
\toprule
\rowcolor{naranja}
\textbf{Modelo} & \textbf{Filas entrenamiento} & \textbf{max\_\allowbreak{}depth} & \textbf{n\_\allowbreak{}estimators} & \textbf{learning\_\allowbreak{}rate} & \textbf{min\_\allowbreak{}child\_\allowbreak{}weight} \\
\midrule
\rowcolor{filaclara} xgboost\_\allowbreak{}residual & 552 & 5 & 300 & 0,01 & 3 \\
xgboost\_\allowbreak{}alone & 889 & 5 & 100 & 0,05 & 10 \\
\bottomrule
\end{tabular}
%
  }
  \caption{\headlinecolor{Hiperparámetros seleccionados por búsqueda en rejilla de 54 puntos (idéntico para ambos modelos XGBoost) sobre un \emph{holdout} interno cronológico.}}
  \label{tab:hiperparametros}
\end{table}

La variable individual con mayor ganancia en este modelo es
\texttt{feat\_\allowbreak{}day\_\allowbreak{}type\_\allowbreak{}festivo} (ganancia 0,125;
tabla \ref{tab:top15_residual}, capítulo \ref{cap:resultados}), lo que apunta a
que la etapa 2 sí aprende, a partir de las variables exógenas, la señal de calendario que
la etapa 1 no puede ver por diseño: el mecanismo que motiva la arquitectura híbrida se
observa de forma medible, con independencia de que el resultado agregado del
capítulo \ref{cap:resultados} termine siendo, o no, favorable a la arquitectura completa.

Conviene dejar constancia, en este punto y no como un descubrimiento tardío en la
discusión de resultados, de una limitación estructural del esquema de \emph{stacking}
OOF empleado. Los cinco modelos LSTM que generan los residuos de entrenamiento se ajustan,
según el fold, sobre entre 337 y 748 filas, mientras que el modelo LSTM final que se
emplea en inferencia, el descrito en la sección \ref{sec:etapa1_lstm}, se entrena sobre
las 889 filas completas y es, previsiblemente, algo mejor que cualquiera de los modelos
parciales que generaron los residuos de entrenamiento. En consecuencia, la etapa 2 aprende
a corregir una distribución de residuo ligeramente más ancha que la que el sistema
realmente enfrenta en inferencia. Este desajuste es inherente a cualquier esquema de
\emph{stacking} fuera de muestra —la alternativa, usar residuos \emph{in-sample} del
modelo final, es estrictamente peor por las razones expuestas en la sección
\ref{sec:esquema_oof}— y se acepta aquí como el coste asumido de evitar una fuga de
información mucho más grave.

\section{Modelos de control}\label{sec:modelos_control}

Para poder atribuir el desempeño del híbrido a su arquitectura, y no simplemente al hecho
de disponer de variables exógenas, es necesario compararlo contra un control que utilice
exactamente las mismas 161 variables sin pasar por ninguna etapa LSTM previa. Este control,
denominado \texttt{xgboost\_alone}, entrena un modelo XGBoost directamente sobre la
demanda total, con idéntico procedimiento de ajuste de hiperparámetros que el modelo
residual de la sección anterior y sobre la misma matriz de características. Su función es
aislar si la etapa 1 aporta valor incremental a la predicción final o si, por el
contrario, introduce ruido que un modelo de una sola etapa evitaría directamente.

Una decisión que podría parecer, a primera vista, una inconsistencia entre ambos modelos
es deliberada: \texttt{xgboost\_alone} se entrena sobre las 889 filas completas del
conjunto de entrenamiento posterior al calentamiento, no sobre las 552 filas del conjunto
residual de la etapa 2. La razón es que \texttt{xgboost\_alone} no depende en ningún
momento de las predicciones OOF de la etapa 1, de modo que ni la exclusión de las 200
filas sin cobertura OOF ni la exclusión del fold 1 degenerado le son aplicables:
imponérselas de todos modos, en nombre de una supuesta igualdad de condiciones, habría
privado al control de un tercio de su historial de entrenamiento disponible por una razón
que le es completamente ajena. La matriz de variables y el procedimiento de ajuste de
hiperparámetros son, por tanto, idénticos entre ambos modelos XGBoost; el tamaño del
conjunto de entrenamiento, legítimamente, no lo es.

La propuesta inicial de este TFM contemplaba, junto a XGBoost, un segundo representante
de la familia de aprendizaje automático tabular (bosques aleatorios, \emph{random forest}). Se ha prescindido
de él de forma deliberada: ambos son métodos de conjunto sobre árboles de decisión y
comparten la misma capacidad de representación sobre esta matriz de características, de
modo que un segundo representante habría añadido una fila más a la comparación maestra
sin poner a prueba ninguna hipótesis distinta. El esfuerzo de control se ha destinado, en
su lugar, a los experimentos de \emph{paridad informativa} de la sección
\ref{sec:paridad_informativa} (ablación de bloques de variables y LSTM con acceso a la
matriz predictora completa), diseñados para discriminar entre las explicaciones rivales del
resultado central de este trabajo.

Como complemento a este control directo, la sección \ref{sec:analisis_sensibilidad} del
capítulo \ref{cap:resultados} introduce una variante adicional de la etapa 2,
\texttt{hybrid\_weighted}, que no constituye un modelo nuevo, sino una reponderación del
mismo modelo residual: idéntica en las 552 filas de entrenamiento, las 161 variables, la
semilla aleatoria y los hiperparámetros (cargados directamente del modelo de la sección
\ref{sec:etapa2_xgboost} en lugar de reajustarse de forma independiente, precisamente
para que ambas variantes sean comparables sin introducir una segunda fuente de
variación), y que difiere únicamente en que asigna un peso de muestra triple a las filas
correspondientes a días de calendario irregular (festivos y días puente), un 8,2\,\% del
conjunto residual. El objetivo de esta variante es comprobar si concentrar la capacidad
del modelo en los días donde la etapa 1 falla con mayor intensidad (festivos y puentes,
según el diagnóstico de la sección \ref{sec:diagnostico_previo}) mejora la corrección precisamente en
esos días; su evaluación completa, con un criterio de éxito fijado antes de observar el
resultado, se desarrolla en la sección \ref{sec:experimento_a} y no se anticipa aquí.

\section{Ensamble por combinación de pronósticos}\label{sec:metodologia_ensemble}

SARIMAX (sección \ref{sec:baselines}) y el control \texttt{xgboost\_alone} de la sección anterior
no solo se emplean por separado como adversarios del híbrido: sus predicciones, ya
generadas y persistidas, permiten formular sin ningún coste de entrenamiento adicional
un tercer candidato —una combinación lineal de ambos— cuya justificación conviene
formalizar aquí, como parte del diseño experimental, y no introducir por primera vez
al presentar sus resultados en el capítulo \ref{cap:resultados}. Dados los pronósticos
puntuales de ambos modelos en el instante $t$, $\hat y_t^{\text{SARIMAX}}$ y
$\hat y_t^{\text{XGB}}$, el ensamble se define como la combinación lineal convexa
\[
  \hat y_t = w_1\, \hat y_t^{\text{SARIMAX}} + w_2\, \hat y_t^{\text{XGB}},
  \qquad w_1 + w_2 = 1,\ w_1, w_2 \geq 0,
\]
una operación puramente aritmética sobre predicciones ya calculadas, sin reentrenar
ningún modelo ni introducir ningún parámetro adicional que deba ajustarse por
descenso de gradiente.

La justificación de que una combinación tan simple pueda superar a ambos componentes
individuales no es heurística: descansa en la identidad de reducción de varianza de Bates y
Granger y en la condición de Granger sobre conjuntos de información distintos, ambas
expuestas en la sección \ref{sec:arquitecturas_hibridas}. Lo que importa fijar aquí es que
SARIMAX y \texttt{xgboost\_alone} \emph{sí} satisfacen esa condición: sus conjuntos de
información se solapan pero no se anidan. El primero condiciona sobre el historial completo
de la serie mediante su componente $(p,d,q)\times(P,D,Q,s)$ con diferenciación estacional
y solo cinco indicadores de calendario; el segundo dispone de una memoria de la serie
limitada a siete retardos y cuatro estadísticos móviles, pero añade los retardos de
operador, los armónicos de Fourier, la codificación completa de calendario y la
meteorología retardada, sin estructura autorregresiva ni diferenciación. El híbrido
secuencial de este TFM, en cambio, no la satisface plenamente, porque el pasado reciente que ve su
etapa 1 está representado en la matriz de su etapa 2 (sección
\ref{sec:resultado_negativo_teoria}).

Se contemplan, por diseño y antes de observar ningún resultado, dos estrategias de
ponderación. La primera, \texttt{ensemble\_equal}, fija $w_1 = w_2 = 0{,}5$: no requiere
ningún dato adicional para determinarse y actúa como el punto de referencia neutro de
la combinación. La segunda, \texttt{ensemble\_inverse\_mae}, pondera cada componente en
proporción inversa a su MAE sobre la partición de validación, de modo que el modelo
históricamente más preciso recibe mayor peso; sus coeficientes se derivan
exclusivamente de la validación y nunca de la partición de test, preservando la
separación estricta entre selección de configuración y evaluación final que gobierna
todo este TFM. Ambas variantes se evalúan en el capítulo \ref{cap:resultados}
(secciones \ref{sec:analisis_sensibilidad} y \ref{sec:mecanismo_ensemble}), donde se
cuantifica empíricamente la correlación de error $\rho$ entre SARIMAX y
\texttt{xgboost\_alone} y se pone la ganancia observada en relación con la reducción de
varianza que la identidad de Bates y Granger anticipa; se adelanta únicamente que ese contraste
sitúa a \texttt{ensemble\_equal} y \texttt{ensemble\_inverse\_mae} como los dos modelos
más precisos de todo el estudio (MAE de test 139.725 y 139.681, respectivamente),
prácticamente indistinguibles entre sí, lo que hace preferible la variante
equiponderada por no depender de ningún dato de validación para fijar sus pesos.
